\documentclass[10pt,a4paper]{amsart}
\usepackage{amsmath,amssymb,amsthm,graphicx}
\usepackage{verbatim}
\newenvironment{reference}{\comment}{\endcomment}
%\newenvironment{reference}{}{}
\newenvironment{slogan}{\comment}{\endcomment}
\newenvironment{history}{\comment}{\endcomment}

% For commutative diagrams we use Xy-pic
\usepackage[all]{xy}

% We use 2cell for 2-commutative diagrams.
\xyoption{2cell}
% Theorem environments.
%
\theoremstyle{plain}
\newtheorem{theorem}[subsection]{Theorem}
\newtheorem{proposition}[subsection]{Proposition}
\newtheorem{lemma}[subsection]{Lemma}

\theoremstyle{definition}
\newtheorem{definition}[subsection]{Definition}
\newtheorem{example}[subsection]{Example}
\newtheorem{exercise}[subsection]{Exercise}
\newtheorem{situation}[subsection]{Situation}

\theoremstyle{remark}
\newtheorem{remark}[subsection]{Remark}
\newtheorem{remarks}[subsection]{Remarks}

\numberwithin{equation}{subsection}

% Macros
%
\def\lim{\mathop{\mathrm{lim}}\nolimits}
\def\colim{\mathop{\mathrm{colim}}\nolimits}
\def\Spec{\mathop{\mathrm{Spec}}}
\def\Hom{\mathop{\mathrm{Hom}}\nolimits}
\def\Ext{\mathop{\mathrm{Ext}}\nolimits}
\def\SheafHom{\mathop{\mathcal{H}\!\mathit{om}}\nolimits}
\def\SheafExt{\mathop{\mathcal{E}\!\mathit{xt}}\nolimits}
\def\Sch{\mathit{Sch}}
\def\Mor{\mathop{\mathrm{Mor}}\nolimits}
\def\Ob{\mathop{\mathrm{Ob}}\nolimits}
\def\Sh{\mathop{\mathit{Sh}}\nolimits}
\def\NL{\mathop{N\!L}\nolimits}
\def\CH{\mathop{\mathrm{CH}}\nolimits}
\def\proetale{{pro\text{-}\acute{e}tale}}
\def\etale{{\acute{e}tale}}
\def\QCoh{\mathit{QCoh}}
\def\Ker{\mathop{\mathrm{Ker}}}
\def\Im{\mathop{\mathrm{Im}}}
\def\Coker{\mathop{\mathrm{Coker}}}
\def\Coim{\mathop{\mathrm{Coim}}}

% Boxtimes
%
\DeclareMathSymbol{\boxtimes}{\mathbin}{AMSa}{"02}

%
% Macros for moduli stacks/spaces
%
\def\QCohstack{\mathcal{QC}\!\mathit{oh}}
\def\Cohstack{\mathcal{C}\!\mathit{oh}}
\def\Spacesstack{\mathcal{S}\!\mathit{paces}}
\def\Quotfunctor{\mathrm{Quot}}
\def\Hilbfunctor{\mathrm{Hilb}}
\def\Curvesstack{\mathcal{C}\!\mathit{urves}}
\def\Polarizedstack{\mathcal{P}\!\mathit{olarized}}
\def\Complexesstack{\mathcal{C}\!\mathit{omplexes}}
% \Pic is the operator that assigns to X its picard group, usage \Pic(X)
% \Picardstack_{X/B} denotes the Picard stack of X over B
% \Picardfunctor_{X/B} denotes the Picard functor of X over B
\def\Pic{\mathop{\mathrm{Pic}}\nolimits}
\def\Picardstack{\mathcal{P}\!\mathit{ic}}
\def\Picardfunctor{\mathrm{Pic}}
\def\Deformationcategory{\mathcal{D}\!\mathit{ef}}
\usepackage{cleveref}
\crefname{originalequation}{original equation}{original equations}
\Crefname{originalequation}{Original equation}{Original equations}
\makeatletter
\newlabel{equation-epsilon}{{08H7}{source}{}{sourcepointer-1}{}}
\newlabel{equation-epsilon@cref}{{[originalequation][0][]08H7}{[1][0][]source}{}{}{}}
\makeatother
\crefname{originalreference}{original reference}{original references}
\Crefname{originalreference}{Original reference}{Original references}
\makeatletter
\newlabel{lemma-derived-quasi-coherent-small-etale-site}{{071Q}{source}{}{sourcepointer-2}{}}
\newlabel{lemma-derived-quasi-coherent-small-etale-site@cref}{{[originalreference][0][]071Q}{[1][0][]source}{}{}{}}
\makeatother
\crefname{originalreference}{original reference}{original references}
\Crefname{originalreference}{Original reference}{Original references}
\makeatletter
\newlabel{lemma-epsilon-flat}{{08H8}{source}{}{sourcepointer-3}{}}
\newlabel{lemma-epsilon-flat@cref}{{[originalreference][0][]08H8}{[1][0][]source}{}{}{}}
\makeatother
\crefname{originallemma}{original lemma}{original lemmas}
\Crefname{originallemma}{Original lemma}{Original lemmas}
\makeatletter
\newlabel{perfect-lemma-tor-qc-qs}{{08EA}{source}{}{sourcepointer-4}{}}
\newlabel{perfect-lemma-tor-qc-qs@cref}{{[originallemma][0][]08EA}{[1][0][]source}{}{}{}}
\makeatother
\crefname{originallemma}{original lemma}{original lemmas}
\Crefname{originallemma}{Original lemma}{Original lemmas}
\makeatletter
\newlabel{spaces-cohomology-lemma-etale-pull-push-split}{{0HAP}{source}{}{sourcepointer-5}{}}
\newlabel{spaces-cohomology-lemma-etale-pull-push-split@cref}{{[originallemma][0][]0HAP}{[1][0][]source}{}{}{}}
\makeatother
\crefname{originallemma}{original lemma}{original lemmas}
\Crefname{originallemma}{Original lemma}{Original lemmas}
\makeatletter
\newlabel{spaces-morphisms-lemma-pushforward}{{03M9}{source}{}{sourcepointer-6}{}}
\newlabel{spaces-morphisms-lemma-pushforward@cref}{{[originallemma][0][]03M9}{[1][0][]source}{}{}{}}
\makeatother
\title{Tor amplitude against all quasi-coherent sheaves on a quasi-separated space}
\author{The Stacks Project Authors}
\date{stacks-project:spaces-perfect.tex:local-snapshot-20261012; GNU FDL 1.2 or later}
\begin{document}
\maketitle
\noindent\textit{Editorial note.} All delivered mathematical text is the author's own, copied from the pinned Stacks Project chapter file spaces-perfect.tex; no statement or proof step was rewritten, reordered or removed.\par\medskip
\noindent\textit{Editorial note.} The counted proof is the author's complete argument: the proof is complete in both directions, closing with the schemes case\par\medskip
\noindent\textit{Editorial note.} The label lemma-tor-qc-qs names the counted statement, and the proof environment immediately following it in the source is delivered verbatim as its proof.\par\medskip
\noindent\textit{Editorial note.} The counted statement is carried by the author's own lemma environment in the chapter source.\par\medskip
\noindent\textit{Allowed standard prerequisites (editorial).}\begin{itemize}
\item tor amplitude of an object of the derived category of modules
\item etale neighbourhoods U to X with affine U, quasi-compact separated by quasi-separatedness, and pushforward preserving quasi-coherence
\item the etale pull-push splitting lemma exhibiting a module as a direct summand of j\textasciicircum{}*j\_*G
\item the comparison of the derived category of quasi-coherent modules on a space with the small etale derived category and flatness of the localization functor
\item the schemes version of the tor-amplitude criterion
\end{itemize}\medskip
\ifdefined\hypertarget\hypertarget{sourcepointer-1}{}\fi
\noindent\textit{Source reference.} Original equation 08H7: sources/stacks-spaces-perfect/pinned-original-source.tex, source lines 320--320. Prerequisite pointer invoked at proof line(s) 3408; it names a standard fact of the same chapter that the counted argument uses, and every step of the counted proof built on it is retained verbatim in the delivered body.\par\medskip
\ifdefined\hypertarget\hypertarget{sourcepointer-2}{}\fi
\noindent\textit{Source reference.} Original reference 071Q: sources/stacks-spaces-perfect/pinned-original-source.tex, source lines 402--402. Prerequisite pointer invoked at proof line(s) 3404; it names a standard fact of the same chapter that the counted argument uses, and every step of the counted proof built on it is retained verbatim in the delivered body.\par\medskip
\ifdefined\hypertarget\hypertarget{sourcepointer-3}{}\fi
\noindent\textit{Source reference.} Original reference 08H8: sources/stacks-spaces-perfect/pinned-original-source.tex, source lines 330--330. Prerequisite pointer invoked at proof line(s) 3416; it names a standard fact of the same chapter that the counted argument uses, and every step of the counted proof built on it is retained verbatim in the delivered body.\par\medskip
\ifdefined\hypertarget\hypertarget{sourcepointer-4}{}\fi
\noindent\textit{Source reference.} Original lemma 08EA: sources/stacks-chow/tags.txt, source lines 10666--10666. Prerequisite pointer invoked at proof line(s) 3422; it names a standard fact of a sibling chapter that the counted argument uses, and every step of the counted proof built on it is retained verbatim in the delivered body.\par\medskip
\ifdefined\hypertarget\hypertarget{sourcepointer-5}{}\fi
\noindent\textit{Source reference.} Original lemma 0HAP: sources/stacks-chow/tags.txt, source lines 21569--21569. Prerequisite pointer invoked at proof line(s) 3394; it names a standard fact of a sibling chapter that the counted argument uses, and every step of the counted proof built on it is retained verbatim in the delivered body.\par\medskip
\ifdefined\hypertarget\hypertarget{sourcepointer-6}{}\fi
\noindent\textit{Source reference.} Original lemma 03M9: sources/stacks-chow/tags.txt, source lines 4741--4741. Prerequisite pointer invoked at proof line(s) 3387; it names a standard fact of a sibling chapter that the counted argument uses, and every step of the counted proof built on it is retained verbatim in the delivered body.\par\medskip
\begin{lemma}
\label{lemma-tor-qc-qs}
Let $S$ be a scheme. Let $X$ be a quasi-separated algebraic space over $S$.
Let $E$ be an object of $D_\QCoh(\mathcal{O}_X)$. Let $a \leq b$.
The following are equivalent
\begin{enumerate}
\item $E$ has tor amplitude in $[a, b]$, and
\item for all $\mathcal{F}$ in $\QCoh(\mathcal{O}_X)$
we have $H^i(E \otimes_{\mathcal{O}_X}^\mathbf{L} \mathcal{F}) = 0$
for $i \not \in [a, b]$.
\end{enumerate}
\end{lemma}

\begin{proof}
It is clear that (1) implies (2). Assume (2). Let $j : U \to X$ be
an \'etale morphism with $U$ affine. As $X$ is quasi-separated, we
see that $j$ is both quasi-compact and separated.
Hence $j_*$ transforms quasi-coherent
modules into quasi-coherent modules (Morphisms of Spaces, Lemma
\ref{spaces-morphisms-lemma-pushforward}).
Take any quasi-coherent module $\mathcal{G}$ on $U$.
By assumption we see that
$H^i(E \otimes_{\mathcal{O}_X}^\mathbf{L} j_*\mathcal{G})$ vanishes
for $i \not \in [a, b]$. Pulling back by the flat morphism $j$ we find that
$H^i(E|_U \otimes_{\mathcal{O}_U}^\mathbf{L} j^*j_*\mathcal{G})$
vanishes for $i \not \in [a, b]$. By
Cohomology of Spaces, Lemma \ref{spaces-cohomology-lemma-etale-pull-push-split}
we see that $\mathcal{G}$ is a direct summand of $j^*j_*\mathcal{G}$.
It follows that condition (2) implies the vanishing of
$H^i(E|_U \otimes_{\mathcal{O}_U}^\mathbf{L} \mathcal{G})$
for $i \not \in [a, b]$ for all quasi-coherent $\mathcal{O}_U$-modules
$\mathcal{G}$. Since it suffices to prove that $E|_U$ has tor amplitude
in $[a, b]$ we reduce to the case where $X$ is representable.

\medskip\noindent
If $X$ is representable by a scheme $X_0$ then
(Lemma \ref{lemma-derived-quasi-coherent-small-etale-site})
we can write $E = \epsilon^*E_0$ where $E_0$ is an object of
$D_\QCoh(\mathcal{O}_{X_0})$ and
$\epsilon : X_\etale \to (X_0)_{Zar}$ is as in
(\ref{equation-epsilon}). For every quasi-coherent module
$\mathcal{F}_0$ on $X_0$ the module $\epsilon^*\mathcal{F}_0$
is quasi-coherent on $X$ and
$$
H^i(E \otimes_{\mathcal{O}_X}^\mathbf{L} \epsilon^*\mathcal{F}_0)
=
\epsilon^*H^i(E_0 \otimes_{\mathcal{O}_{X_0}}^\mathbf{L} \mathcal{F}_0)
$$
as $\epsilon$ is flat (Lemma \ref{lemma-epsilon-flat}).
Moreover, the vanishing of these sheaves for $i \not \in [a, b]$
implies the same thing for
$H^i(E_0 \otimes_{\mathcal{O}_{X_0}}^\mathbf{L} \mathcal{F}_0)$
by the same lemma. Thus we've reduced the problem to the case
of schemes which is treated in
Derived Categories of Schemes, Lemma \ref{perfect-lemma-tor-qc-qs}.
\end{proof}

\end{document}
